\section{Preliminaries}
\noindent
Let $R=\bigoplus_{i\geq 0}R_i$ be a standard graded $\K$-algebra. A \emph{Koszul filtration} of $R$ is a family of ideals $\mathcal{F}$ satisfying:
\begin{enumerate}
    \item[\tagKFlin] Every ideal in $\mathcal{F}$ is generated by elements of $R_1$.
    \item[\tagKFmax] The maximal ideal $\m_R$ belongs to $\mathcal{F}$.
    \item[\tagKFcol] For every nonzero ideal $I\in \mathcal{F}$, there exists an ideal $J\in \mathcal{F}$ properly contained in $I$ such that $I/J$ is cyclic and $J:I\in \mathcal{F}$.
\end{enumerate}
Recall that for ideals $I,J \in R$, the colon ideal is $J:I=\{f\in R \mid fI \subseteq J\}$. We will use the notation $J:x$ for $J:(J+(x))$.

Note that by (KF3), $\mathcal{F}$ contains a chain of ideals from the zero ideal $(0)$ to the maximal ideal $\m_R$: 
$$(0)\subset I_1 \subset I_2\subset\dots\subset\m_R$$
\noindent
where $I_i$ is minimally generated by $i$ generators and such that each colon ideal $I_i:I_{i+1}$ is linear and again in $\mathcal{F}$. This will be important for the proof of Theorem A. For every ideal $I$ in a Koszul filtration, $R/I$ has a linear minimal free resolution \cite{conca2012koszulalgebrasregularity}.

Let $A=\K [x_1,...,x_n]$ be a polynomial algebra, and let $S\subset A$ be any subalgebra. The following Lemma describes the colon ideal $(f):g$ where $(f)$ and $(f,g)$ are both ideals in a subalgebra $S\subset A$. This description resembles other descriptions of colon ideals, such as Theorem $14$ in Section $4$, Chapter $4$ of Cox, Little, and O'Shea \cite{cox2015ideals}.

\begin{lem}\label{coli}
 Let $f$ and $g$ be two elements in the subalgebra $S\subset A$, where $f=f_0h$ and $g=g_0h$ such that $\gcd_A(f_0,g_0)=1$. We can describe the colon ideal $(f):_Sg$ as follows:
    $$(f):_Sg=\{f_0p\mid f_0p, g_0p \in S \text{ and } p\in A\}$$
\end{lem}
\begin{proof}
We first show that the containment $\subseteq$ holds. Recall that by the definition of a colon ideal, $a\in (f):_Sg$ implies that 
$$ag=bf \Longleftrightarrow ag_0h=bf_0h\Longleftrightarrow ag_0=bf_0$$

\noindent
for some $b\in S$. Since we are in a unique factorisation domain, this implies $f_0|a$ and $g_0|b$. Hence we can express $a=f_0p_1$ and $b=g_0p_2$.
$$f_0p_1g_0=g_0p_2f_0\Longleftrightarrow p_1=p_2$$
\noindent
Thus $a=f_0p_1$ belongs to the right-hand side. For the inverse containment $\supseteq$, take any 
$ c \in \{f_0p\mid f_0p, g_0p \in S \text{ and } p\in A\}$. Then $cg=(f_0p)(g_0h)=(f_0h)(g_0p) \in (f)$, hence $c\in (f):_Sg$.

    
\end{proof}
